% ----- CHAUPAI 25 -----
\subsection*{\deva{पञ्चविंशतितम चौपाई} (Twenty-Fifth Chaupai)}
\addcontentsline{toc}{subsection}{\deva{पञ्चविंशतितम चौपाई} (Twenty-Fifth Chaupai) — \deva{नासै रोग...}}

\begin{minipage}[t]{0.55\textwidth}
\vspace{0pt}

{\large \linespread{1.5}\selectfont
\deva{नासै रोग हरै सब पीरा।}\\
\deva{जपत निरन्तर हनुमत बीरा॥}
\par}

\vspace{0.5em}

{\large \linespread{1.5}\selectfont
\telu{నాసై రోగ హరై సబ పీరా।}\\
\telu{జపత నిరంతర హనుమత బీరా॥}
\par}

\vspace{0.5em}

{\large \linespread{1.3}\selectfont
\textit{nāsai roga harai saba pīrā |}\\
\textit{japata nirantara hanumata bīrā ||}
\par}
\end{minipage}%
\hfill
\begin{minipage}[t]{0.42\textwidth}
\vspace{0pt}
\begin{tcolorbox}[anchorbox]
\textbf{Consistency is the Cure:} Just as studying requires daily effort, maintaining mental health and overcoming life's pains (\textit{Peera}) requires constant (\textit{Nirantara}), disciplined practice. You cannot cram for character; it must be built consistently over time.
\vspace{0.5em}\newline He proved that consistent, relentless action acts as the ultimate medicine for healing deep physical and mental exhaustion.\newline\null\hfill\href{https://www.valmiki.iitk.ac.in/sloka?field_kanda_tid=6&language=dv&field_sarga_value=74}{\textit{Yuddha Kanda, 74:73-75}}
\end{tcolorbox}
\end{minipage}

\vspace{0.5em}
\subsubsection*{\textbf{Visual Rhythm Map:}}
\begin{table}[h!]
\centering
\renewcommand{\arraystretch}{1.2}
\setlength{\tabcolsep}{0pt}
\begin{tabularx}{\textwidth}{|*{16}{>{\centering\arraybackslash\hsize=1.0\hsize}X|}}
\hline
\multicolumn{4}{|>{\centering\arraybackslash\hsize=4.0\hsize}X|}{\textbf{\guru{ना}\guru{सै}} \par (4)} &
\multicolumn{3}{>{\centering\arraybackslash\hsize=3.0\hsize}X|}{\textbf{\guru{रो}\laghu{ग}} \par (3)} &
\multicolumn{3}{>{\centering\arraybackslash\hsize=3.0\hsize}X|}{\textbf{\laghu{ह}\guru{रै}} \par (3)} &
\multicolumn{2}{>{\centering\arraybackslash\hsize=2.0\hsize}X|}{\textbf{\laghu{स}\laghu{ब}} \par (2)} &
\multicolumn{4}{>{\centering\arraybackslash\hsize=4.0\hsize}X|}{\textbf{\guru{पी}\guru{रा}} \par (4)} \\
\hline
\end{tabularx}
\vspace{-1.5em}
\end{table}

\begin{table}[h!]
\centering
\renewcommand{\arraystretch}{1.2}
\setlength{\tabcolsep}{0pt}
\begin{tabularx}{\textwidth}{|*{16}{>{\centering\arraybackslash\hsize=1.0\hsize}X|}}
\hline
\multicolumn{3}{|>{\centering\arraybackslash\hsize=3.0\hsize}X|}{\textbf{\laghu{ज}\laghu{प}\laghu{त}} \par (3)} &
\multicolumn{5}{>{\centering\arraybackslash\hsize=5.0\hsize}X|}{\textbf{\laghu{नि}\guru{रं}\laghu{त}\laghu{र}} \par (5)} &
\multicolumn{4}{>{\centering\arraybackslash\hsize=4.0\hsize}X|}{\textbf{\laghu{ह}\laghu{नु}\laghu{म}\laghu{त}} \par (4)} &
\multicolumn{4}{>{\centering\arraybackslash\hsize=4.0\hsize}X|}{\textbf{\guru{बी}\guru{रा}} \par (4)} \\
\hline
\end{tabularx}
\vspace{-1.5em}
\end{table}

\vspace{1em}
\textbf{ \deva{सरल-संस्कृतम्}:}\\
 \deva{सर्वे रोगाः नश्यन्ति, सर्वाः पीडाः च हरिताः भवन्ति}\\
 \deva{यः निरन्तरं हे वीर हनुमन्! भवतः नाम जपति॥}\\


All diseases are destroyed, and all suffering/pain is taken away\\
for one who continuously chants your name, O brave Hanuman!


\vspace{1em}
\newpage
\textbf{\deva{प्रतिपदार्थः} (Meaning of each word):}
\begin{table}[h!]
\centering
\renewcommand{\arraystretch}{1.2}
\begin{tabularx}{\textwidth}{@{} l l X X @{}}
\toprule
\textbf{Awadhi} & \textbf{Sanskrit} & \textbf{English} & \textbf{Grammar \& Notes} \\
\midrule
\deva{नासै} \gotchaicon / \telu{నాసై} / \textit{nāsai}  &  \deva{नश्यति}  &  Is destroyed  & Verb ending in `-ai` \\
\deva{रोग} / \telu{రోగ} / \textit{roga}  &  \deva{रोगाः}  &  Diseases  &   \\
\deva{हरै} \gotchaicon / \telu{హరై} / \textit{harai}  &  \deva{हरति}  &  Takes away / destroys  & Verb ending in `-ai` \\
\deva{सब पीरा} / \telu{సబ పీరా} / \textit{saba pīrā}  &  \deva{सर्वाः पीडाः}  &  All pain  & \deva{पीडा $\rightarrow$ पीरा} \\ 
\deva{जपत} / \telu{జపత} / \textit{japata}  &  \deva{जपन्}  &  Chanting  & Continuous action `-ta` \\
\deva{निरन्तर} / \telu{నిరంతర} / \textit{nirantara}  &  \deva{निरन्तरम्}  &  Continuously  &   \\
\deva{हनुमत} / \telu{హనుమత} / \textit{hanumata}  &  \deva{हनुमन्}  &  Hanuman  &   \\
\deva{बीरा} / \telu{బీరా} / \textit{bīrā}  &  \deva{वीरः}  &  Brave/Hero  & \par  \deva{} \\ 
\bottomrule
\end{tabularx}
\end{table}

\begin{tcolorbox}[gotchabox]
\begin{itemize}
    \item \textbf{The \deva{ड} $\rightarrow$ \deva{र} Shift:} The Sanskrit word for pain is \deva{पीडा} (Peedaa). Over time, the retroflex 'D' softens into a rolling 'R' sound, giving the beautiful Awadhi word \deva{पीरा} (Peeraa).
    \item \textbf{Verbs ending in -ai:} Just like \deva{गावैं} and \deva{लहै}, the verbs \deva{नासै} (Destroy) and \deva{हरै} (Remove) are standard present-tense conjugations in Awadhi poetry.
\end{itemize}
\end{tcolorbox}

\newpage
